\documentclass[11pt,a4paper]{article}
\usepackage[margin=2.5cm]{geometry}
\usepackage[utf8]{inputenc}
\usepackage[T1]{fontenc}
\usepackage[spanish]{babel}
\usepackage{amsmath,amssymb,amsthm}
\usepackage{xcolor}
\usepackage[most]{tcolorbox}
\usepackage{hyperref}
\usepackage{enumitem}
\usepackage{booktabs}
\usepackage{lmodern}
\usepackage{microtype}

% Colors
\definecolor{itbaBlue}{HTML}{003366}
\definecolor{itbaCyan}{HTML}{6EC6D9}
\definecolor{tipGreen}{HTML}{2E7D32}
\definecolor{tipGreenLight}{HTML}{E8F5E9}
\definecolor{consignaBlue}{HTML}{E3F2FD}
\definecolor{consignaBlueDark}{HTML}{1565C0}
\definecolor{warnOrange}{HTML}{E65100}
\definecolor{warnOrangeLight}{HTML}{FFF3E0}
\definecolor{theoremPurple}{HTML}{4A148C}
\definecolor{theoremPurpleLight}{HTML}{F3E5F5}

\tcbuselibrary{skins,breakable}

\newtcolorbox{consignabox}{
  enhanced, breakable,
  colback=consignaBlue, colframe=consignaBlueDark,
  boxrule=1.5pt, arc=4pt,
  title={\textbf{Consigna}},
  fonttitle=\bfseries\color{white},
  attach boxed title to top left={yshift=-2mm, xshift=4mm},
  boxed title style={colback=consignaBlueDark, arc=2pt, boxrule=0pt}
}

\newtcolorbox{tipsbox}{
  enhanced, breakable,
  colback=tipGreenLight, colframe=tipGreen,
  boxrule=1.5pt, arc=4pt,
  title={\textbf{Tips y Anotaciones}},
  fonttitle=\bfseries\color{white},
  attach boxed title to top left={yshift=-2mm, xshift=4mm},
  boxed title style={colback=tipGreen, arc=2pt, boxrule=0pt}
}

\newtcolorbox{teoriabox}[1]{
  enhanced, breakable,
  colback=theoremPurpleLight, colframe=theoremPurple,
  boxrule=1.5pt, arc=4pt,
  title={\textbf{#1}},
  fonttitle=\bfseries\color{white},
  attach boxed title to top left={yshift=-2mm, xshift=4mm},
  boxed title style={colback=theoremPurple, arc=2pt, boxrule=0pt}
}

\newtcolorbox{warnbox}[1]{
  enhanced, breakable,
  colback=warnOrangeLight, colframe=warnOrange,
  boxrule=1.5pt, arc=4pt,
  title={\textbf{#1}},
  fonttitle=\bfseries\color{white},
  attach boxed title to top left={yshift=-2mm, xshift=4mm},
  boxed title style={colback=warnOrange, arc=2pt, boxrule=0pt}
}

\hypersetup{
  colorlinks=true,
  linkcolor=itbaBlue,
  urlcolor=itbaCyan,
  pdftitle={Ejercicios RSA -- ITBA Criptografía},
  pdfauthor={ITBA Criptografía y Seguridad}
}

\title{
  \vspace{-1cm}
  {\color{itbaBlue}\rule{\linewidth}{2pt}}\\[0.3cm]
  {\Huge\bfseries\color{itbaBlue} RSA: Clave Pública y Firma Digital}\\[0.2cm]
  {\large\color{itbaCyan} Ejercicios de Parciales -- ITBA Criptografía y Seguridad (72.44)}\\[0.2cm]
  {\color{itbaBlue}\rule{\linewidth}{2pt}}
}
\author{}
\date{\today}

\begin{document}
\maketitle
\tableofcontents
\newpage

%=============================================================
\section{Teoría: RSA}
%=============================================================

\subsection{Generación de claves (GenRSA)}

\begin{teoriabox}{GenRSA --- Generación de claves RSA}
\begin{enumerate}
  \item Elegir $p, q$ primos. Calcular $n = p \cdot q$ (el ``módulo'').
  \item $\varphi(n) = (p-1)(q-1)$ (función de Euler).
  \item Elegir $e$ tal que $1 < e < \varphi(n)$ y $\gcd(e, \varphi(n)) = 1$ (coprimo).
  \item Calcular $d$ tal que $e \cdot d \equiv 1 \pmod{\varphi(n)}$ (inverso modular, Euclides extendido).
  \item \textbf{Clave pública:} $pk = \langle n, e \rangle$. \quad \textbf{Clave privada:} $sk = \langle n, d \rangle$.
\end{enumerate}
\[
  \text{Cifrar: } c = m^e \bmod n, \qquad \text{Descifrar: } m = c^d \bmod n.
\]
\end{teoriabox}

\subsection{¿Por qué funciona? (Teorema de Euler)}

\begin{teoriabox}{Corrección de RSA (Euler/Fermat)}
Como $d \equiv e^{-1} \pmod{\varphi(n)}$, existe $k$ entero tal que $ed = 1 + k\varphi(n)$. Entonces:
\[
  c^d \bmod n = (m^e)^d \bmod n = m^{ed} \bmod n = m^{1+k\varphi(n)} \bmod n = m \cdot (m^{\varphi(n)})^k \bmod n.
\]
Por el \textbf{teorema de Euler}: $m^{\varphi(n)} \equiv 1 \pmod{n}$ (cuando $\gcd(m,n)=1$). Luego $c^d \equiv m \pmod{n}$. Cifrado y descifrado son funciones inversas con el par de claves correcto.
\end{teoriabox}

\subsection{Problemas del Textbook RSA}

\begin{warnbox}{Problemas del Textbook RSA (sin padding)}
\begin{itemize}
  \item \textbf{Es determinístico:} el mismo $m$ con el mismo $pk$ siempre da el mismo $c$. No es CPA-Secure (el atacante puede probar mensajes).
  \item \textbf{$e$ y $m$ pequeños:} si $m^e < n$, no hay reducción modular y se puede calcular $\sqrt[e]{c}$ directamente (raíz $e$-ésima entera).
  \item \textbf{Módulos repetidos:} si dos pares comparten el mismo $n$, pueden factorizarlo.
  \item \textbf{Ataque de texto plano elegido (CCA):} Textbook RSA es homomórfico ($E(m_1)\cdot E(m_2) = E(m_1\cdot m_2)$), lo que permite ataques sofisticados.
\end{itemize}
\end{warnbox}

\subsection{PKCS \#1 v1.5: Padding aleatorio}

\begin{teoriabox}{PKCS \#1 v1.5 --- Making RSA Non-Deterministic}
Sea $k$ la longitud de $n$ en bytes. Se construye el mensaje con padding:
\[
  m' = \texttt{00}\,\|\,\texttt{02}\,\|\,r\,\|\,\texttt{00}\,\|\,m
\]
donde $r$ son bytes aleatorios $\ne 0$ (longitud $k - |m| - 3$). Luego se cifra $m'$.
\begin{itemize}
  \item \textbf{El padding aporta no-determinismo:} cada cifrado es diferente aunque el mensaje sea el mismo.
  \item \textbf{Se cree que es CPA-Secure}, pero no es CCA-Secure (existen ataques como Bleichenbacher).
  \item En la práctica, \textbf{siempre usar padding} (PKCS o OAEP). Nunca ``textbook RSA''.
\end{itemize}
\end{teoriabox}

\subsection{Firma digital RSA}

\begin{teoriabox}{RSA como esquema de firma}
En cifrado: la clave \textbf{pública} cifra, la \textbf{privada} descifra.

En firma: la clave \textbf{privada} firma, la \textbf{pública} verifica. Las claves tienen roles diferentes.

\medskip
\textbf{Firma estándar (Textbook RSA sign):}
\begin{itemize}
  \item $\text{Sign}_{sk}(m) = m^d \bmod n = s$ (la firma).
  \item $\text{Vrfy}_{pk}(m, s) = [s^e \bmod n \stackrel{?}{=} m]$.
\end{itemize}

\textbf{Firma con hash (RSA-FDH / Hashed RSA):}
\begin{itemize}
  \item $\text{Sign}_{sk}(m) = [H(m)]^d \bmod n = s$.
  \item $\text{Vrfy}_{pk}(m, s)$: calcular $s^e \bmod n$ y verificar que $= H(m)$.
  \item \textbf{Por qué hashear:} (1) RSA solo puede firmar mensajes del tamaño del módulo; (2) el hash evita ataques existenciales (falsificación sin conocer el mensaje).
\end{itemize}
\end{teoriabox}

\subsection{PKI y Certificados}

\begin{teoriabox}{Infraestructura de Clave Pública (PKI)}
El problema de la clave pública: ¿cómo sé que la clave pública que recibo realmente pertenece a quien dice ser?

\textbf{Solución:} \textbf{Certificados digitales} y \textbf{Autoridades de Certificación (CA)}.
\begin{itemize}
  \item Una CA firma un certificado que asocia identidad $\leftrightarrow$ clave pública.
  \item El certificado contiene: nombre del titular, clave pública, fecha de vencimiento, firma de la CA.
  \item Para confiar en un certificado, hay que confiar en la CA que lo firmó (cadena de confianza).
\end{itemize}
\end{teoriabox}

\subsection{Resumen de roles de las claves}

\begin{center}
\renewcommand{\arraystretch}{1.4}
\begin{tabular}{lll}
\toprule
\textbf{Operación} & \textbf{Clave que actúa} & \textbf{Clave que verifica/deshace} \\
\midrule
Cifrado & Pública (del destinatario) & Privada (del destinatario) \\
Firma digital & Privada (del firmante) & Pública (del firmante) \\
\bottomrule
\end{tabular}
\end{center}

\begin{warnbox}{Trampa clásica de parcial}
``Los sistemas de cifrado de clave pública utilizan una clave para cifrar y otra para firmar.'' \textbf{FALSO.} En cifrado: la pública cifra, la privada descifra. En firma: la privada firma, la pública verifica. Las claves no están asignadas a una operación específica (cifrar vs. firmar) sino a un \emph{rol} (quién las usa). Esta distinción es clave en los parciales.
\end{warnbox}

%=============================================================
\section{Parcial 2015 -- Ejercicio 4 (Primera parte)}
\label{sec:rsa2015a}
%=============================================================

\begin{consignabox}
En un esquema de clave pública, seleccionar la opción correcta:
\begin{enumerate}[label=(\alph*)]
  \item La clave pública se mantiene en secreto y la clave privada se distribuye autenticada.
  \item La clave privada se mantiene en secreto y la clave pública se distribuye autenticada.
  \item La clave privada se mantiene en secreto y la clave pública también.
  \item La clave privada se mantiene en secreto y la clave pública se mantiene en un KDC.
\end{enumerate}
\end{consignabox}

\paragraph{Resolución.}

\textbf{Respuesta: (b).}

En un esquema de clave pública (asimétrico):
\begin{itemize}
  \item La \textbf{clave privada} se mantiene \textbf{en secreto} --- solo el dueño la conoce.
  \item La \textbf{clave pública} se \textbf{distribuye} libremente, pero debe ser \textbf{autenticada} (para evitar ataques de suplantación de identidad). Esto es exactamente el rol de los certificados digitales y la PKI.
\end{itemize}

\emph{Por qué las demás son falsas:}
\begin{itemize}
  \item \textbf{(a)} Invierte los roles: la pública es la que se distribuye.
  \item \textbf{(c)} Ambas en secreto sería un sistema simétrico, no asimétrico.
  \item \textbf{(d)} La clave pública en un KDC (Key Distribution Center) es el modelo de Kerberos (clave simétrica), no de clave pública.
\end{itemize}

\begin{tipsbox}
\textbf{Patrón:} cuando el enunciado pregunte sobre la distribución de claves en criptografía asimétrica, la respuesta es siempre: \textbf{privada = secreta, pública = distribuida autenticada}. La palabra ``autenticada'' es clave: si distribuyo la clave pública sin autenticar, un atacante puede hacer un ataque MitM suplantando la identidad.
\end{tipsbox}

%=============================================================
\section{Parcial 2015 -- Ejercicio 4 (Segunda parte)}
\label{sec:rsa2015b}
%=============================================================

\begin{consignabox}
Para reducir ciertos ataques en el esquema de firma RSA se sugiere usar $\pi = (\text{Gen}, \text{Sign}, \text{Vrfy})$ donde $s$ (la firma) se calcula como $s := [H(m)]^d \bmod N$, siendo $(N,d) = sk$ y $H$ una función de hash resistente a colisiones. Para verificar la firma, ¿cuál es el proceso correcto que debe seguir Vrfy?

\begin{enumerate}[label=(\alph*)]
  \item Vrfy debe recibir $s$ y efectuar $s^e \bmod N$ y $H(s^e \bmod N)$.
  \item Vrfy debe recibir $s$ y efectuar $s^{(-1)^e} \bmod N$.
  \item Vrfy debe recibir $m$ y $s$ y efectuar $s^e \bmod N$ y comparar con $H(m)$.
  \item Vrfy debe recibir $m$ y $s$ y efectuar $m^e \bmod N$ y comparar con $H(s)$.
\end{enumerate}
\end{consignabox}

\paragraph{Resolución.}

\textbf{Respuesta: (c).}

El proceso es:
\begin{enumerate}
  \item \textbf{Firma (Sign):} $s = [H(m)]^d \bmod N$.
  \item \textbf{Verificación (Vrfy):} dado $(m, s)$:
  \begin{itemize}
    \item Calcular $s^e \bmod N$ --- esto ``deshace'' la firma con la clave pública.
    \item Calcular $H(m)$ --- el hash del mensaje recibido.
    \item Verificar que $s^e \bmod N \stackrel{?}{=} H(m)$.
  \end{itemize}
\end{enumerate}

\emph{Por qué las demás son falsas:}
\begin{itemize}
  \item \textbf{(a)} Solo recibe $s$, no $m$: no puede calcular $H(m)$ para comparar.
  \item \textbf{(b)} $s^{(-1)^e}$ no tiene sentido algebraico en este contexto.
  \item \textbf{(d)} $m^e$ cifra el mensaje original, no está relacionado con la firma; $H(s)$ tampoco tiene sentido.
\end{itemize}

\begin{tipsbox}
\textbf{Mnemotecnia para firma RSA con hash:}

\medskip
``\textbf{Sign priva-da, Hash primero; Vrfy publi-ca, descifra y compara con Hash del mensaje.}''

\medskip
\begin{itemize}
  \item \textbf{Sign:} toma $m$ $\to$ hashea $\to$ eleva a la $d$ (privada) $\to$ produce $s$.
  \item \textbf{Vrfy:} toma $(m, s)$ $\to$ eleva $s$ a la $e$ (pública) $\to$ compara con $H(m)$.
  \item \textbf{Key insight:} $s^e = [H(m)^d]^e = H(m)^{de} = H(m)^{1+k\varphi(n)} = H(m)$ (por Euler). Por eso la comparación funciona.
  \item \textbf{Por qué hashear y no firmar $m$ directamente:} (1) RSA solo firma mensajes $< n$; (2) sin hash, el adversario puede fabricar firmas válidas para mensajes que no eligió (ataque existencial).
\end{itemize}
\end{tipsbox}

%=============================================================
\section{Parcial 2015 -- Ejercicio 5}
\label{sec:rsa2015c}
%=============================================================

\begin{consignabox}
\textbf{Protocolo de autenticación RSA.} El siguiente protocolo de autenticación mutua usa certificados y firmas RSA:
\begin{enumerate}
  \item $A \to B$: $n_A$ (nonce de A)
  \item $B \to A$: $\text{Cert}_B$, $n_B$, $\text{Sign}_B(n_A, n_B)$
  \item $A \to B$: $\text{Cert}_A$, $\text{Sign}_A(n_B, n_A)$
\end{enumerate}
\textbf{Mostrar cómo un adversario $C$ puede autenticarse como $A$ ante $B$.}
\end{consignabox}

\paragraph{Resolución.}

Este protocolo tiene una vulnerabilidad de \textbf{ataque de reflexión} (o ataque de entrelazado). El adversario $C$ puede autenticarse como $A$ ante $B$ sin conocer la clave privada de $A$:

\textbf{Sesión 1 (C ataca a B fingiendo ser A):}
\begin{enumerate}
  \item $C \to B$: $n_C$ (C inicia, fingiendo ser A, con un nonce propio)
  \item $B \to C$: $\text{Cert}_B$, $n_B$, $\text{Sign}_B(n_C, n_B)$
\end{enumerate}
Ahora $C$ necesita $\text{Sign}_A(n_B, n_C)$ para completar el paso 3. $C$ abre una \textbf{sesión paralela}:

\textbf{Sesión 2 (C ataca a B abriendo otra sesión):}
\begin{enumerate}
  \item $C \to B$: \textbf{$n_B$} (C envía el nonce de B como su nonce inicial)
  \item $B \to C$: $\text{Cert}_B$, $n_{B'}$, $\text{Sign}_B(n_B, n_{B'})$
\end{enumerate}
Pero $C$ necesita que B le devuelva $\text{Sign}_A(n_B, n_C)$\ldots\ El ataque es posible si $A=B$ (sesión de reflexión) o si $B$ es convencido de firmar lo necesario.

\textbf{En la variante más directa (C se hace pasar por A ante B):}
\begin{itemize}
  \item $C$ inicia sesión con $B$ enviando $n_A$ (copiando o eligiendo un nonce).
  \item $B$ responde con $\text{Sign}_B(n_A, n_B)$: el desafío es $n_B$.
  \item $C$ necesita $\text{Sign}_A(n_B, n_A)$. Si $C$ puede inducir a $A$ a firmar eso (por ejemplo, abriendo una sesión con $A$ donde $A$ firme el nonce $n_B$), $C$ puede completar la autenticación.
  \item El protocolo \textbf{no incluye el nombre del destinatario dentro de lo firmado}, por lo que la misma firma puede usarse en diferentes contextos.
\end{itemize}

\textbf{Defensa:} incluir la identidad del destinatario en el mensaje firmado:
\begin{itemize}
  \item Paso 2: $\text{Sign}_B(A, n_A, n_B)$ (incluir el nombre del interlocutor).
  \item Paso 3: $\text{Sign}_A(B, n_B, n_A)$ (idem).
\end{itemize}
Así una firma obtenida en una sesión no puede reutilizarse en otra.

\begin{tipsbox}
\textbf{Análisis de protocolos de autenticación --- checklist:}
\begin{itemize}
  \item \textbf{¿Los mensajes firmados incluyen la identidad del interlocutor?} Si no, son vulnerables a ataque de reflexión.
  \item \textbf{¿Los nonces son suficientemente aleatorios?} Si no, hay ataques de replay.
  \item \textbf{¿El protocolo autentica ambas partes?} Un protocolo unilateral no da autenticación mutua.
  \item \textbf{¿Los certificados están verificados?} Si el certificado no se valida, hay MitM.
  \item \textbf{Nonce = Number used once:} aleatorio, para demostrar ``freshness'' (que el mensaje es reciente, no un replay).
  \item La vulnerabilidad de reflexión es un clásico: el adversario abre dos sesiones en paralelo y usa la firma de una para completar la otra.
\end{itemize}
\end{tipsbox}

%=============================================================
\section{Parcial 2025 Q1 -- Ejercicio 5b}
\label{sec:rsa2025Q1}
%=============================================================

\begin{consignabox}
\textbf{Verdadero o Falso (justificar):} ``Los sistemas de encripción basados en clave pública utilizan una de las claves para encriptar/desencriptar y la otra para firmar.''
\end{consignabox}

\paragraph{Resolución.}

\textbf{FALSO.}

La afirmación mezcla incorrectamente las operaciones. La correcta distinción es:

\begin{center}
\renewcommand{\arraystretch}{1.4}
\begin{tabular}{lll}
\toprule
\textbf{Operación} & \textbf{Usa clave} & \textbf{Verifica/deshace con} \\
\midrule
Cifrado & Pública (del destinatario) & Privada (del destinatario) \\
Descifrado & Privada (del destinatario) & -- \\
Firma & Privada (del firmante) & -- \\
Verificación de firma & -- & Pública (del firmante) \\
\bottomrule
\end{tabular}
\end{center}

No hay ``una clave para cifrar/descifrar'' y ``otra para firmar''. Las claves tienen roles determinados por su naturaleza (pública vs. privada):
\begin{itemize}
  \item La \textbf{clave privada} se usa para: descifrar mensajes dirigidos a uno mismo, y firmar mensajes propios.
  \item La \textbf{clave pública} se usa para: cifrar mensajes dirigidos al dueño, y verificar firmas del dueño.
\end{itemize}

\begin{tipsbox}
\textbf{Este es un ejercicio de comprensión conceptual pura.} El error del enunciado es pensar en las claves como ``una para X y otra para Y'', cuando en realidad cada clave aparece en dos operaciones distintas (una como agente, otra como verificador). Recordar: \emph{la clave que ``actúa'' siempre es la privada (descifrar, firmar); la que ``verifica'' siempre es la pública (cifrar para alguien, verificar su firma)}.
\end{tipsbox}

%=============================================================
\section{Parcial 2025 Q2 -- Ejercicio 5c}
\label{sec:rsa2025Q2}
%=============================================================

\begin{consignabox}
\textbf{Verdadero o Falso (justificar):} ``El uso de padding aleatorio en la implementación del algoritmo de clave pública de RSA es para que el algoritmo sea seguro a ataque de textos cifrados elegidos.''
\end{consignabox}

\paragraph{Resolución.}

\textbf{FALSO.}

El padding aleatorio (PKCS \#1 v1.5 y OAEP) convierte al RSA determinístico en \textbf{no-determinístico}, lo que lo hace \textbf{CPA-Secure} (seguro ante ataque de texto plano elegido), \textbf{no} CCA-Secure (seguro ante ataque de texto cifrado elegido).

La distinción:
\begin{itemize}
  \item \textbf{CPA (Chosen Plaintext Attack):} el atacante elige plaintexts y ve sus cifrados. El RSA textbook (determinístico) falla porque el mismo $m$ siempre da el mismo $c$. El padding aleatorio hace que el mismo $m$ dé distintos $c$ $\Rightarrow$ CPA-Secure.
  \item \textbf{CCA (Chosen Ciphertext Attack):} el atacante elige ciphertexts y ve sus descifrados. PKCS \#1 v1.5 \textbf{no} protege contra CCA: existen ataques como el de \textbf{Bleichenbacher (1998)} que explotan el oráculo de descifrado para recuperar el texto plano.
  \item Para CCA-Security se necesita \textbf{OAEP} (Optimal Asymmetric Encryption Padding) u otras construcciones probablemente seguras bajo CCA.
\end{itemize}

\begin{tipsbox}
\textbf{Jerarquía de seguridad que hay que recordar:}

\medskip
\noindent CCA-Secure $\Rightarrow$ CPA-Secure $\Rightarrow$ Secure against eavesdropping

\medskip
\begin{itemize}
  \item \textbf{RSA Textbook (sin padding):} solo seguro si el atacante no puede ver ningún cifrado (poco útil en la práctica).
  \item \textbf{RSA + PKCS \#1 v1.5:} CPA-Secure (no determinístico), pero Bleichenbacher mostró que NO es CCA-Secure.
  \item \textbf{RSA + OAEP:} se demuestra CCA-Secure bajo suposiciones criptográficas estándar.
  \item \textbf{Trampa del enunciado:} la afirmación dice ``CCA'' pero el padding solo da ``CPA''. Esta confusión entre CPA y CCA es uno de los errores más penalizados en parciales de criptografía.
\end{itemize}
\end{tipsbox}

%=============================================================
\section{Resumen y Tabla de Ejercicios RSA}
%=============================================================

\begin{center}
\renewcommand{\arraystretch}{1.4}
\begin{tabular}{llll}
\toprule
\textbf{Año} & \textbf{Ej.} & \textbf{Tema} & \textbf{Respuesta} \\
\midrule
2015 & 4a & Distribución de claves en PKI & (b): privada secreta, pública autenticada \\
2015 & 4b & Verificación de firma RSA con hash & (c): Vrfy recibe $(m,s)$, calcula $s^e$ y $H(m)$ \\
2015 & 5 & Ataque a protocolo de autenticación RSA & Ataque de reflexión / entrelazado \\
2025 Q1 & 5b & Roles de las claves (V/F) & FALSO: claves no se asignan a operaciones \\
2025 Q2 & 5c & Padding y CPA/CCA (V/F) & FALSO: padding da CPA-Secure, no CCA-Secure \\
\bottomrule
\end{tabular}
\end{center}

\begin{warnbox}{Trampas frecuentes en RSA}
\begin{enumerate}
  \item \textbf{CPA vs CCA:} el padding aleatorio hace RSA CPA-Secure. Para CCA-Secure se necesita OAEP u otras construcciones.
  \item \textbf{Roles de las claves:} la privada descifra y firma; la pública cifra y verifica. No hay ``una para cifrar/descifrar y otra para firmar''.
  \item \textbf{Firma RSA hashed:} Sign calcula $[H(m)]^d$; Vrfy recibe $(m,s)$ y compara $s^e$ con $H(m)$.
  \item \textbf{Verificación requiere el mensaje $m$:} en hashed RSA, Vrfy no puede funcionar solo con $s$ (necesita $m$ para calcular $H(m)$).
  \item \textbf{Protocolo de autenticación:} si lo firmado no incluye la identidad del destinatario, hay vulnerabilidad de reflexión.
  \item \textbf{$e$ y $d$ se usan en módulo $\varphi(n)$:} $e \cdot d \equiv 1 \pmod{\varphi(n)}$, no módulo $n$.
\end{enumerate}
\end{warnbox}

\end{document}
